\documentclass[11pt,a4paper]{article}
\usepackage[margin=2.5cm]{geometry}
\usepackage[T1]{fontenc}
\usepackage{lmodern}
\usepackage[hidelinks]{hyperref}
\setlength{\parskip}{0.6em}
\setlength{\parindent}{0pt}
\pagestyle{empty}

\begin{document}

\begin{flushright}
\textbf{Giona Granchelli}\\
+31~(0)6~427433~16\\
giona.granchelli@gmail.com\\
\href{https://gionag.com}{gionag.com}
\end{flushright}

\vspace{1.5em}

Dear Hiring Team,

I am applying for the Engineering Chapter Lead, AI \& Agentic Platforms position because it sits at the intersection of the things I most enjoy doing: building technology that can make a real difference, solving difficult engineering problems, developing people, and continuously learning myself.

I currently work as an Engineering Chapter Lead and Software Engineer IV at ABN AMRO, where I lead a Java and Vue chapter of approximately 45 engineers across five teams. Alongside the people and engineering leadership side of my role, I have become increasingly focused on how AI and agentic systems can improve the way engineering teams work.

As Copilot Champion and AI Adoption Lead, I have been building and introducing governed agentic workflows for software development, incident investigation, migrations, and other engineering activities. One example was an agentic migration approach we used across ten legacy applications, reducing what we estimated would have required more than 18 months of engineering effort to approximately four months. Just as importantly, I have worked on the governance around these systems, including approval controls, auditable execution, and human review.

Outside my role at ABN AMRO, I continue to explore the engineering challenges behind AI platforms in depth. I created TramAI, an open-source JVM runtime for governed AI applications and multi-agent workflows, and I experiment extensively with RAG, model evaluation, local inference, fine-tuning, and AI orchestration. I enjoy understanding how these systems work rather than only consuming them as tools.

What attracts me to this opportunity at ING is the chance to take that work to a much larger level. I want to help build AI capabilities that are genuinely useful, reliable, secure, and capable of improving how people work across an organization. Doing that in a large regulated environment, where architecture, governance and engineering quality matter as much as experimentation, is exactly the kind of challenge I am looking for.

I also care strongly about the people side of engineering. One of the parts of being a Chapter Lead that I value most is helping engineers grow, sharing what I learn, creating good engineering practices, and enabling teams to solve problems better together. I want my next role to keep challenging me technically while also giving me the opportunity to help others become stronger engineers.

ING is a company I have respected for a long time for its engineering culture, scale, and position within European banking. The opportunity to contribute to its AI direction, learn from people operating at that scale, and take on problems that will push me further is one I would be very proud to pursue.

I would welcome the opportunity to discuss what I could bring to the chapter and what we could build together.

Kind regards,

\textbf{Giona Granchelli}

\end{document}
